\section{Introduction}\label{sec:intro}

Federated Learning (FL) trains a shared model on data that never leaves the
devices that produced it~\cite{mcmahan2017fedavg}. On smartphones, every
participation is paid in battery: a selected phone downloads the global model,
trains it on its local data, and uploads the update over the cellular network.
Battery life is a primary concern of smartphone users~\cite{ferreira2011battery},
and a phone that runs out of charge leaves the federation, together with its
data.

\textbf{State of the art.}
Energy-aware FL mostly acts at the server, by choosing \emph{who} trains.
SAGE~\cite{savoia2026sage} ranks clients by state of charge (SoC) and data
divergence and selects them by integer programming; ESCS~\cite{maciel2024escs}
filters clients by residual battery and network quality and weighs them by
their loss. Once selected, however, clients run the same local work, or a
reduced amount decided by the server: SAGE, for instance, halves the epochs of
selected clients with low energy. FedProx~\cite{li2020fedprox} tolerates
heterogeneous local work through a proximal term, but as a nuisance to
correct, not as a knob to save energy.

\textbf{Motivation.}
How much local work a phone should do in a round depends on two things: how
much work it \emph{can afford}, given its charge, the load of the other
applications and whether it is plugged in, and how much work is \emph{worth
doing}, given how poorly the current global model fits its data. Both are
known exactly only on the device, and only when training starts. The server
sees them through reports that are old by construction, since a client is
contacted only once every few rounds. Refreshing this view is not free:
polling every device at each round costs a radio session per phone, and
sending the model to a phone that then declines wastes a full download. With
a fixed local workload, many phones run out of battery before training is
over (Sec.~\ref{sec:results}).

\textbf{Approach.}
We propose SAGE-smart, a client-driven assignment of the local work. The server
keeps the decision it can take well: who participates, from a ranking that
balances SoC, data divergence and the time since the last participation.
Instead of a fixed workload, it sends an \emph{envelope}: a range of epochs, a
set of batch sizes, a round deadline and a battery reserve. The client chooses
its epochs and batch size inside the envelope when training starts (late
binding), using its fresh state and the loss of the global model on a sample
of its data. A lightweight proposal, sent before the model, lets a phone that
cannot afford even the minimum work decline at the cost of a control message
instead of a model download.

\textbf{Contributions.} In detail, our main contributions are:
\begin{itemize}
  \item A client-driven local work assignment for FL on battery-powered
  devices: the server selects the participants and proposes an envelope, and
  each client binds its epochs and batch size at training time, within its
  battery reserve and the round deadline (Sec.~\ref{sec:methodology}).
  \item A characterization, in messages and joules, of three ways to bring
  the client state into this decision: model-first negotiation, lightweight
  negotiation, and state broadcast with a server-side decision. Sending the
  proposal before the model makes a refusal an order of magnitude cheaper
  than a wasted download, with no fixed cost per round
  (Sec.~\ref{sec:protocols}).
  \item An energy-faithful FL simulator built on Flower~\cite{beutel2020flower},
  based on the speed and power of real smartphones~\cite{xu2022mandheling},
  three battery models, the LTE radio~\cite{huang2012lte,narayanan2021variegated},
  a background workload and the recharge habits of real
  users~\cite{ferreira2011battery}. Every joule spent by a device, from
  training to idle, is charged to its battery (Sec.~\ref{sec:system}).
  % TODO mention the code release, if we release it
  \item An evaluation with ResNet-18 on CIFAR-100 against FedAvg, FedProx,
  SAGE and ESCS, under different battery models and levels of data
  heterogeneity. With respect to FedAvg, SAGE-smart \rev{uses 40\% less energy
  and depletes 8 phones instead of 58, with 2.8 more accuracy points}; with
  respect to SAGE and ESCS, on average, it \rev{gains 9.7 accuracy points with
  39\% less energy} (Sec.~\ref{sec:results}).
  % TODO add the workload stress test (when the negotiation pays off)
\end{itemize}
% NOTE: results from campaign v2 (3 seeds, datasheet battery, mean of the three
% beta, SAGE-smart without FedNova). SAGE and ESCS = mean of sage_nm and the
% four ESCS variants (sd, sp, md, mp), datasheet battery.
% NOTE: FedAvg with B = 16 reaches a higher final accuracy (report 12); the
% claim on accuracy must be stated against the B = 128 baselines, or dropped.
